\documentclass[11pt]{article}

\usepackage[T1]{fontenc}
\usepackage[utf8]{inputenc}
\usepackage{lmodern}
\usepackage{microtype}
\usepackage[a4paper,margin=1in]{geometry}
\usepackage{amsmath,amssymb,amsthm,mathtools}
\usepackage{booktabs,longtable,array}
\usepackage{enumitem}
\usepackage{xcolor}
\usepackage{hyperref}
\usepackage[nameinlink,noabbrev]{cleveref}
\usepackage{fancyhdr}
\usepackage{url}

\hypersetup{
  colorlinks=true,
  linkcolor=blue!55!black,
  citecolor=green!40!black,
  urlcolor=magenta!55!black,
  pdftitle={Digit-transfer separation and synchronized gap ladders for the Alaoglu--Erdos conjecture},
  pdfauthor={OpenAI GPT-5.6 Pro},
  pdfsubject={Research progress report on the Alaoglu--Erdos conjecture}
}

\pagestyle{fancy}
\fancyhf{}
\lhead{Alaoglu--Erd\H{o}s progress report}
\rhead{Digit transfer and synchronized gaps}
\cfoot{\thepage}
\setlength{\headheight}{14pt}
\allowdisplaybreaks
\setlist[itemize]{topsep=3pt,itemsep=2pt,parsep=1pt}
\setlist[enumerate]{topsep=3pt,itemsep=2pt,parsep=1pt}

\newtheorem{theorem}{Theorem}[section]
\newtheorem{proposition}[theorem]{Proposition}
\newtheorem{lemma}[theorem]{Lemma}
\newtheorem{corollary}[theorem]{Corollary}
\newtheorem{conjecture}[theorem]{Conjecture}
\newtheorem{problem}[theorem]{Problem}
\newtheorem{criterion}[theorem]{Criterion}
\theoremstyle{definition}
\newtheorem{definition}[theorem]{Definition}
\newtheorem{remark}[theorem]{Remark}
\newtheorem{example}[theorem]{Example}
\newtheorem{warning}[theorem]{Warning}

\newcommand{\Z}{\mathbb Z}
\newcommand{\N}{\mathbb N}
\newcommand{\Q}{\mathbb Q}
\newcommand{\R}{\mathbb R}
\newcommand{\C}{\mathbb C}
\newcommand{\Qp}{\mathbb Q_p}
\newcommand{\Zp}{\mathbb Z_p}
\newcommand{\vTwo}{v_2}
\newcommand{\vThree}{v_3}
\newcommand{\ord}{\operatorname{ord}}
\newcommand{\dist}{\operatorname{dist}}
\newcommand{\rad}{\operatorname{rad}}
\newcommand{\sgn}{\operatorname{sgn}}
\newcommand{\Frac}{\operatorname{Frac}}
\newcommand{\PhiDT}{\Phi}
\newcommand{\PsiDT}{\Psi}
\newcommand{\e}{\mathrm e}

\title{\textbf{Digit-Transfer Separation, Macroscopic Defect Laws,\\
and a Determinant-One Gap Ladder}\\[0.6em]
\large New progress toward the Alaoglu--Erd\H{o}s conjecture}
\author{OpenAI GPT-5.6 Pro}
\date{22 August 2026}

\begin{document}
\maketitle

\begin{center}
\fbox{\begin{minipage}{0.92\textwidth}
\textbf{Status declaration.}
This report does \emph{not} claim a proof of the Alaoglu--Erd\H{o}s conjecture.
It develops a new, self-contained digit-transfer framework, proves several
rigorous separation and defect theorems, constructs an exact synchronized
continued-fraction ladder, and isolates two sharply formulated missing
arithmetic inputs.  Statements labeled as theorems or propositions below
are proved; proposed closing steps are explicitly labeled as problems or
criteria.
\end{minipage}}
\end{center}

\begin{abstract}
Let
\[
  \vartheta=\log_2 3,\qquad \varrho=\vartheta^{-1}=\log_3 2.
\]
The Alaoglu--Erd\H{o}s conjecture asks whether integers
\(M=2^x\) and \(A=3^x\) can occur for a nonintegral real \(x\).
The supplied unified report already develops the standard transcendence
barrier, primitive normalization, pure-radical reduction, determinant and
Pad\'e constructions, local methods, and numerous approximation-theoretic
routes.  The present report deliberately avoids reproducing that material.

Our first new device is the exact binary-to-ternary digit-transfer map
\[
  \Phi\!\left(\sum \epsilon_j2^j\right)=\sum \epsilon_j3^j
\]
and its ternary-to-binary retraction
\[
  \Psi\!\left(\sum \delta_j3^j\right)=\sum \delta_j2^j.
\]
We prove their carry calculus, sharp global power bounds, summatory and
Dirichlet-series identities, and a two-prime-adic retraction theorem.
For a hypothetical pair \(A=M^\vartheta\), the defects
\[
 D_k=A^k-\Phi(M^k),\qquad E_k=\Psi(A^k)-M^k
\]
are positive integers.  They satisfy exact interval-valuation identities,
\(0<E_k\le D_k\), and \(E_k\le 2D_k^\varrho\).

The strongest structural result is a \emph{prefix-separation theorem}:
when the base-2 expansion of \(M^k\) and the base-3 expansion of \(A^k\)
are aligned at their common most-significant position, their common prefix
is either empty or exactly a leading \(1\) followed by zeros.  At the first
mismatch, the ternary digit is strictly larger.  Combining this with an
explicit lower bound for a linear form in \(\log M\) and \(\log 2\) places
the first mismatch within \(O(\log k)\) places of the leading digit.  We
also prove positive-density and effective all-\(k\) macroscopic lower bounds
for both defects.

A second new construction starts from the pure-radical normalization in the
supplied report.  Continued-fraction convergents to a single logarithm
produce simultaneous near-power gaps at bases \(2\) and \(3\) with the same
small logarithmic error and adjacent exponent determinant exactly one.  On
one sign subsequence, digit transfer gives an exact ``gap-descent square''.
The square is not contractive: we quantify the scale mismatch that prevents
it from closing the conjecture.  This no-go calculation leaves a precise
research target---a prime-transfer or carry-rigidity theorem strong enough
to connect the two synchronized gaps.
\end{abstract}

\tableofcontents
\newpage

\section{Scope, imported facts, and nonduplication}

\subsection{The conjecture and hypothetical notation}

The conjecture is
\begin{equation}\label{eq:AE}
  x\in\R,\quad 2^x\in\Z,\quad 3^x\in\Z
  \quad\Longrightarrow\quad x\in\Z.
\end{equation}
Write
\begin{equation}\label{eq:MA-beta}
  M=2^x,\qquad A=3^x,
  \qquad \vartheta=\log_2 3,
  \qquad \varrho=\log_3 2.
\end{equation}
Then
\begin{equation}\label{eq:power-relation}
  A=M^\vartheta,\qquad M=A^\varrho,
  \qquad \log_2 M=\log_3 A=x.
\end{equation}
A nonintegral rational \(x\) is impossible by unique factorization, so a
hypothetical counterexample has irrational \(x\).  Standard transcendence
theorems moreover force it beyond the algebraic-irrational case.  These
reductions, and their relationship to the four-exponentials frontier, are
part of the supplied report and are not re-proved here.

Throughout the main digit-transfer analysis we therefore make the following
standing hypothesis solely for contradiction-seeking purposes.

\begin{definition}[Hypothetical pair]\label{def:hyp-pair}
A \emph{hypothetical pair} is a pair of integers \(M,A>1\) such that
\(A=M^\vartheta\), with neither \(M\) a power of \(2\) nor \(A\) a power of
\(3\).  We set
\[
  \beta=\log_2 M=\log_3 A\notin\Q.
\]
\end{definition}

\subsection{What is imported from the supplied report}

Only the following established reductions are imported.

\begin{enumerate}
\item A counterexample, if one exists, may be organized through a least
positive generator and a primitive normalization.
\item In one convenient branch there are integers \(a,d\ge0\), an odd
integer \(w>1\), and an integer \(A>1\) such that
\begin{equation}\label{eq:radical-import}
  M=2^a w^d,
  \qquad \eta=w^\vartheta,
  \qquad \eta^d=\frac{A}{3^a},
\end{equation}
and \(\eta\) has exact algebraic degree \(d\).
\item Generic adjacent near-power gaps admit a determinant-controlled gcd
bound: after removing primes dividing the two bases, the gcd of two gaps is
controlled by the determinant of their exponent vectors.  We invoke this
only in \cref{sec:gap-ladder}, where our new exponent vectors have
determinant exactly one.
\end{enumerate}

No determinant expansions, Pad\'e approximants, finite-difference
constructions, automatic-sequence surveys, or local analytic arguments from
the supplied manuscript are repeated.  The maps \(\Phi\) and \(\Psi\), the
prefix-separation theorem, the defect interval identities, and the
synchronized determinant-one ladder below are new organizational devices
relative to that manuscript.

\subsection{Why digit transfer is a natural test}

The equality \(A^k=(M^k)^\vartheta\) converts each binary atom exactly into
a ternary atom:
\[
  (2^j)^\vartheta=3^j.
\]
For a binary expansion \(M^k=\sum\epsilon_j2^j\), convexity gives
\[
  A^k=\left(\sum\epsilon_j2^j\right)^\vartheta
  >\sum\epsilon_j3^j
\]
unless there is only one nonzero bit.  The right side is an integer with
ternary digits \(0,1\).  Conversely, concavity of \(t^\varrho\) compares the
ternary expansion of \(A^k\) with \(M^k\).  This creates two integral defects
whose arithmetic can be studied exactly, rather than only estimated by
real analysis.

\section{The two digit-transfer maps}
\label{sec:maps}

\begin{definition}[Finite digit-transfer maps]\label{def:maps}
For \(n\in\N_0\), write its canonical binary and ternary expansions as
\[
 n=\sum_{j\ge0}\epsilon_j(n)2^j,
 \quad \epsilon_j(n)\in\{0,1\},
 \qquad
 n=\sum_{j\ge0}\delta_j(n)3^j,
 \quad \delta_j(n)\in\{0,1,2\}.
\]
Define
\begin{equation}\label{eq:PhiPsi-def}
 \Phi(n)=\sum_{j\ge0}\epsilon_j(n)3^j,
 \qquad
 \Psi(n)=\sum_{j\ge0}\delta_j(n)2^j.
\end{equation}
\end{definition}

Thus \(\Phi\) reinterprets a binary word as a ternary word, while \(\Psi\)
reinterprets a ternary word as a binary digit string that is then evaluated
without first normalizing digits \(2\).

\subsection{Radix recurrences}

\begin{proposition}[Exact recurrences]\label{prop:recurrences}
For all \(n\ge0\),
\begin{align}
 \Phi(2n)&=3\Phi(n),
 &\Phi(2n+1)&=3\Phi(n)+1,\label{eq:Phi-rec}\\
 \Psi(3n+r)&=2\Psi(n)+r
 &&(r=0,1,2).\label{eq:Psi-rec}
\end{align}
\end{proposition}

\begin{proof}
Appending a binary digit \(r\in\{0,1\}\) multiplies the transferred ternary
word by \(3\) and adds \(r\).  Appending a ternary digit
\(r\in\{0,1,2\}\) multiplies its binary evaluation by \(2\) and adds \(r\).
\end{proof}

The increment sequences are unexpectedly simple.

\begin{theorem}[Valuation increments]\label{thm:increments}
For every integer \(n\ge1\),
\begin{align}
 \Phi(n)-\Phi(n-1)
   &=\frac{3^{v_2(n)}+1}{2},\label{eq:Phi-inc}\\
 \Psi(n)-\Psi(n-1)
   &=2-2^{v_3(n)}.\label{eq:Psi-inc}
\end{align}
\end{theorem}

\begin{proof}
Let \(r=v_2(n)\).  Passing from \(n-1\) to \(n\) changes the trailing binary
word \(0\,1^r\) to \(1\,0^r\).  Under \(\Phi\), the increment is
\[
  3^r-\sum_{j=0}^{r-1}3^j
  =3^r-\frac{3^r-1}{2}
  =\frac{3^r+1}{2}.
\]
For \(r=v_3(n)\), the trailing ternary word changes from
\(0\,2^r\) to \(1\,0^r\).  Its binary evaluation changes by
\[
  2^r-2\sum_{j=0}^{r-1}2^j
  =2^r-2(2^r-1)
  =2-2^r.
\]
\end{proof}

\begin{corollary}[Summatory valuation identities]\label{cor:summatory}
For \(N\ge1\),
\begin{align}
 \Phi(N)&=\frac{N}{2}+\frac12\sum_{n\le N}3^{v_2(n)},
 \label{eq:Phi-sum}\\
 \Psi(N)&=2N-\sum_{n\le N}2^{v_3(n)}.
 \label{eq:Psi-sum}
\end{align}
\end{corollary}

\begin{proof}
Sum \cref{eq:Phi-inc,eq:Psi-inc} from \(1\) to \(N\), using
\(\Phi(0)=\Psi(0)=0\).
\end{proof}

\subsection{Dirichlet series and log-periodic scaling}

\begin{proposition}[Valuation Dirichlet series]\label{prop:dirichlet}
For a prime \(p\), a positive real \(c\), and
\(\Re(s)>\max\{1,\log_p c\}\),
\begin{equation}\label{eq:generic-vp-DS}
 \sum_{n\ge1}\frac{c^{v_p(n)}}{n^s}
 =\zeta(s)\frac{1-p^{-s}}{1-cp^{-s}}.
\end{equation}
In particular,
\begin{align}
 \sum_{n\ge1}\frac{3^{v_2(n)}}{n^s}
  &=\zeta(s)\frac{1-2^{-s}}{1-3\,2^{-s}},
  &&\Re(s)>\vartheta,\label{eq:DS-2}\\
 \sum_{n\ge1}\frac{2^{v_3(n)}}{n^s}
  &=\zeta(s)\frac{1-3^{-s}}{1-2\,3^{-s}},
  &&\Re(s)>1.\label{eq:DS-3}
\end{align}
\end{proposition}

\begin{proof}
Write \(n=p^jm\) with \(p\nmid m\).  Then
\[
 \sum_{n\ge1}\frac{c^{v_p(n)}}{n^s}
 =\sum_{j\ge0}(cp^{-s})^j
   \sum_{p\nmid m}\frac1{m^s}
 =\frac{1}{1-cp^{-s}}(1-p^{-s})\zeta(s).
\]
\end{proof}

The extra pole lattice of \cref{eq:DS-2} on
\(\Re(s)=\vartheta\) reflects the exact scaling
\(\Phi(2n)=3\Phi(n)\).  In \cref{eq:DS-3}, the pole of \(\zeta(s)\) at
\(s=1\) gives the linear term \(2N\) in \cref{eq:Psi-sum}; subtracting it
reveals the lower-order \(N^\varrho\)-scale fractal \(\Psi(N)\).

\subsection{Carry monotonicity}

\begin{theorem}[Carry inequalities]\label{thm:carry-ineq}
For all nonnegative integers \(m,n\),
\begin{align}
 \Phi(m+n)&\ge\Phi(m)+\Phi(n),
 &\Phi(mn)&\ge\Phi(m)\Phi(n),\label{eq:Phi-super}\\
 \Psi(m+n)&\le\Psi(m)+\Psi(n),
 &\Psi(mn)&\le\Psi(m)\Psi(n).\label{eq:Psi-sub}
\end{align}
Moreover, \(\Phi\) is strictly increasing.
\end{theorem}

\begin{proof}
Before carrying, addition or multiplication of binary digit strings produces
nonnegative integer coefficients at the weights \(2^j\).  A binary carry
replaces \(2\cdot2^j\) by \(2^{j+1}\).  Under evaluation at \(3\), the same
operation replaces \(2\cdot3^j\) by \(3^{j+1}\), increasing the value by
\(3^j\).  Normalizing all carries therefore proves the two inequalities for
\(\Phi\).  Strict monotonicity also follows from \cref{eq:Phi-inc}.

For ternary strings, a carry replaces \(3\cdot3^j\) by \(3^{j+1}\).  Under
evaluation at \(2\), it replaces \(3\cdot2^j\) by \(2^{j+1}\), decreasing
the value by \(2^j\).  Normalizing all ternary carries proves both
inequalities for \(\Psi\).
\end{proof}

\begin{warning}
The map \(\Psi\) is not order-preserving.  Indeed, its increment in
\cref{eq:Psi-inc} is negative whenever \(v_3(n)\ge2\).  Any argument that
applies monotonicity to \(\Psi\) is therefore invalid.
\end{warning}

\subsection{Composition and a two-prime-adic retraction}

\begin{theorem}[Retraction identities]\label{thm:retraction}
For every \(n\ge0\),
\begin{equation}\label{eq:PsiPhi}
  \Psi(\Phi(n))=n.
\end{equation}
For every \(a\ge0\),
\begin{equation}\label{eq:PhiPsi-major}
  \Phi(\Psi(a))\ge a,
\end{equation}
with equality if and only if the ternary expansion of \(a\) uses only the
digits \(0\) and \(1\).  Consequently
\begin{equation}\label{eq:H-idem}
 H=\Phi\circ\Psi\quad\text{satisfies}\quad H^2=H.
\end{equation}
\end{theorem}

\begin{proof}
The ternary digits of \(\Phi(n)\) are exactly the binary digits of \(n\),
which proves \cref{eq:PsiPhi}.  To compute \(\Phi(\Psi(a))\), first regard
the ternary digits \(\delta_j(a)\in\{0,1,2\}\) as coefficients at binary
weights and then normalize binary carries.  Every carry increases evaluation
at \(3\), by the proof of \cref{thm:carry-ineq}.  Equality occurs precisely
when no digit \(2\) and no subsequent carry is present, namely when all
ternary digits lie in \(\{0,1\}\).  Finally,
\(H^2=\Phi\Psi\Phi\Psi=\Phi\Psi=H\).
\end{proof}

The finite maps extend naturally to incompatible local fields.

\begin{theorem}[Two-prime-adic extension]\label{thm:padic-extension}
The formula
\[
 \Phi\!\left(\sum_{j\ge0}\epsilon_j2^j\right)
   =\sum_{j\ge0}\epsilon_j3^j
\]
defines a continuous injection \(\Phi:\Z_2\to\Z_3\), whose image is the
3-adic Cantor set
\[
 \mathcal C_3=\left\{\sum_{j\ge0}\epsilon_j3^j:
          \epsilon_j\in\{0,1\}\right\}.
\]
The formula
\[
 \Psi\!\left(\sum_{j\ge0}\delta_j3^j\right)
   =\sum_{j\ge0}\delta_j2^j
\]
defines a continuous surjection \(\Psi:\Z_3\to\Z_2\), and
\(\Psi\circ\Phi=\operatorname{id}_{\Z_2}\).  If \(x,y\in\Z_2\), then
\begin{equation}\label{eq:padic-isometry}
  v_3(\Phi(x)-\Phi(y))=v_2(x-y).
\end{equation}
If \(a\equiv b\pmod{3^r}\), then
\begin{equation}\label{eq:Psi-cont}
  \Psi(a)\equiv\Psi(b)\pmod{2^r}.
\end{equation}
\end{theorem}

\begin{proof}
Both series converge in their target local fields.  The first differing
binary digit of \(x,y\) has coefficient \(\pm1\), so it remains the first
nonzero ternary coefficient in the difference; this proves
\cref{eq:padic-isometry} and injectivity.  Surjectivity and the retraction
identity follow by choosing ternary digits in \(\{0,1\}\).  Congruence
\cref{eq:Psi-cont} follows because the first \(r\) ternary digits determine
the output modulo \(2^r\).
\end{proof}

\section{Sharp archimedean power bounds}
\label{sec:power-bounds}

The elementary bounds below are optimal in their constants and will be
used repeatedly.

\begin{theorem}[Sharp global bounds]\label{thm:sharp-bounds}
For every integer \(n\ge1\),
\begin{equation}\label{eq:sharp-Phi}
  \frac12 n^\vartheta\le\Phi(n)\le n^\vartheta,
\end{equation}
and
\begin{equation}\label{eq:sharp-Psi}
  n^\varrho\le\Psi(n)\le2n^\varrho.
\end{equation}
Equality in the upper bound for \(\Phi\), or in the lower bound for \(\Psi\),
occurs exactly when \(n\) is a power of the relevant source base
(\(2\) for \(\Phi\), \(3\) for \(\Psi\)).  The constants \(1/2\) and \(2\)
are best possible.
\end{theorem}

\begin{proof}
Write \(n=\sum\epsilon_j2^j\).  Since \(\vartheta>1\),
\[
 n^\vartheta
 =\left(\sum\epsilon_j2^j\right)^\vartheta
 \ge\sum\epsilon_j(2^j)^\vartheta
 =\Phi(n),
\]
with equality only when a single bit is nonzero.

For the lower bound, use induction on the binary length.  Write
\(n=2^L+r\), \(0\le r<2^L\), and set \(t=r/2^L\).  Assuming
\(\Phi(r)\ge r^\vartheta/2\),
\[
 \Phi(n)=3^L+\Phi(r)
 \ge3^L\left(1+\frac12t^\vartheta\right).
\]
The function
\(2+t^\vartheta-(1+t)^\vartheta\) is decreasing on \([0,1]\) and vanishes
at \(1\), because \(2^\vartheta=3\).  Hence
\[
 1+\frac12t^\vartheta\ge\frac12(1+t)^\vartheta,
\]
which proves the induction.

For \(\Psi\), write \(n=\sum\delta_j3^j\).  Since \(0<\varrho<1\),
subadditivity gives
\[
 n^\varrho\le\sum_j(\delta_j3^j)^\varrho
 =\sum_j\delta_j^\varrho2^j
 \le\sum_j\delta_j2^j=\Psi(n).
\]
Equality requires exactly one digit, equal to \(1\).

For the upper bound, induct on ternary length.  Write
\(n=d3^L+r\), where \(d\in\{1,2\}\), \(0\le r<3^L\), and
\(t=r/3^L\).  The induction hypothesis gives
\[
 \Psi(n)=d2^L+\Psi(r)\le2^L(d+2t^\varrho).
\]
For \(d=1,2\), the function
\[
 2(d+t)^\varrho-d-2t^\varrho
\]
is decreasing on \([0,1]\), and its value at \(1\) is nonnegative
(the \(d=2\) value is exactly \(0\), since \(3^\varrho=2\)).  Therefore
\(d+2t^\varrho\le2(d+t)^\varrho\), completing the induction.

Finally, at \(n=2^L-1\),
\(\Phi(n)=(3^L-1)/2\), so the ratio in the lower bound tends to \(1/2\).
At \(n=3^L-1\), \(\Psi(n)=2(2^L-1)\), and the ratio in the upper bound
tends to \(2\).
\end{proof}

\begin{corollary}[Root limits]\label{cor:root-limits}
For every integer \(m\ge2\),
\begin{align}
 \lim_{k\to\infty}\Phi(m^k)^{1/k}&=m^\vartheta,\label{eq:Phi-root}\\
 \lim_{k\to\infty}\Psi(m^k)^{1/k}&=m^\varrho.\label{eq:Psi-root}
\end{align}
\end{corollary}

\begin{proof}
Take \(k\)-th roots in \cref{eq:sharp-Phi,eq:sharp-Psi}.
\end{proof}

\subsection{Quantitative one-step convexity and concavity}

\begin{proposition}[Binary leading-atom defect]\label{prop:binary-defect}
If \(n=2^L+r\) with \(1\le r<2^L\), then
\begin{align}
 n^\vartheta-\Phi(n)
 &\ge3^L\left[
     \left(1+\frac r{2^L}\right)^\vartheta
     -1-\left(\frac r{2^L}\right)^\vartheta
   \right]\label{eq:binary-defect-exact}\\
 &\ge(\vartheta-1)3^L\frac r{2^L}
 \ge(\vartheta-1)\left(\frac32\right)^L.
 \label{eq:binary-defect-linear}
\end{align}
\end{proposition}

\begin{proof}
Use \(\Phi(r)\le r^\vartheta\).  For
\(f(t)=(1+t)^\vartheta-1-t^\vartheta\), one has
\[
 f'(t)=\vartheta\big((1+t)^{\vartheta-1}-t^{\vartheta-1}\big)
 \ge \frac\vartheta2>\vartheta-1
 \qquad(0\le t\le1),
\]
where the minimum occurs at \(t=1\) and
\(2^{\vartheta-1}=3/2\).  Thus \(f(t)\ge(\vartheta-1)t\).
\end{proof}

\begin{proposition}[Ternary leading-atom defect]\label{prop:ternary-defect}
If \(n=d3^L+r\), where \(d\in\{1,2\}\) and \(0\le r<3^L\), then
\begin{equation}\label{eq:ternary-defect-exact}
 \Psi(n)-n^\varrho
 \ge2^L\left[d+
       \left(\frac r{3^L}\right)^\varrho
       -\left(d+\frac r{3^L}\right)^\varrho\right].
\end{equation}
For \(d=1\),
\begin{equation}\label{eq:ternary-defect-d1}
 \Psi(n)-n^\varrho
 \ge(2-2^\varrho)2^L\left(\frac r{3^L}\right)^\varrho.
\end{equation}
For \(d=2\),
\begin{equation}\label{eq:ternary-defect-d2}
 \Psi(n)-n^\varrho\ge(2-2^\varrho)2^L.
\end{equation}
\end{proposition}

\begin{proof}
Use \(\Psi(r)\ge r^\varrho\).  For \(d=1\), the function
\[
 g(t)=\frac{(1+t)^\varrho-1}{t^\varrho}
\]
is increasing on \((0,1]\), because the sign of its derivative is the sign
of \(1-(1+t)^{\varrho-1}>0\).  Hence
\(g(t)\le g(1)=2^\varrho-1\), proving
\cref{eq:ternary-defect-d1}.  For \(d=2\), the expression
\(2+t^\varrho-(2+t)^\varrho\) is increasing in \(t\) and at \(t=0\)
equals \(2-2^\varrho\).
\end{proof}

\section{The integral defect system}
\label{sec:defects}

Assume henceforth that \((M,A)\) is a hypothetical pair in the sense of
\cref{def:hyp-pair}.  Define, for \(k\ge1\),
\begin{equation}\label{eq:defect-defs}
 C_k=\Phi(M^k),\qquad
 D_k=A^k-C_k,
 \qquad
 E_k=\Psi(A^k)-M^k,
 \qquad
 K_k=\Phi(\Psi(A^k))-A^k.
\end{equation}

\begin{proposition}[Positivity]\label{prop:defect-positive}
For every \(k\ge1\),
\begin{equation}\label{eq:defect-signs}
  D_k>0,
  \qquad E_k>0,
  \qquad K_k\ge0.
\end{equation}
\end{proposition}

\begin{proof}
The integer \(M^k\) is not a power of \(2\), for otherwise \(M\) would be a
power of \(2\).  Strictness in the convexity proof of
\cref{thm:sharp-bounds} gives
\(\Phi(M^k)<(M^k)^\vartheta=A^k\).  Similarly, \(A^k\) is not a power of
\(3\), so \(\Psi(A^k)>(A^k)^\varrho=M^k\).  The assertion for \(K_k\) is
\cref{eq:PhiPsi-major}.
\end{proof}

\subsection{Exact interval-valuation identities}

\begin{theorem}[Defect interval identities]\label{thm:defect-interval}
For every \(k\ge1\),
\begin{align}
 E_k
 &=2D_k-
   \sum_{C_k<n\le A^k}2^{v_3(n)},
 \label{eq:E-interval}\\
 D_k+K_k
 &=\frac12\left(
 E_k+
 \sum_{M^k<n\le M^k+E_k}3^{v_2(n)}
 \right).
 \label{eq:DK-interval}
\end{align}
Consequently,
\begin{equation}\label{eq:E-le-D}
 0<E_k\le D_k
\end{equation}
and
\begin{equation}\label{eq:E-PsiD}
 E_k\le\Psi(D_k)\le2D_k^\varrho.
\end{equation}
\end{theorem}

\begin{proof}
Because \(\Psi(C_k)=\Psi(\Phi(M^k))=M^k\), telescoping
\cref{eq:Psi-inc} over \((C_k,A^k]\) gives
\[
 E_k=\Psi(A^k)-\Psi(C_k)
 =\sum_{C_k<n\le A^k}(2-2^{v_3(n)}),
\]
which is \cref{eq:E-interval}.  There are \(D_k\) terms in the sum and each
\(2^{v_3(n)}\ge1\), proving \(E_k\le D_k\).

Next, \(\Psi(A^k)=M^k+E_k\), so telescoping \cref{eq:Phi-inc} gives
\[
 \Phi(\Psi(A^k))-\Phi(M^k)
 =\frac12\sum_{M^k<n\le M^k+E_k}(3^{v_2(n)}+1).
\]
The left side is \((A^k+K_k)-(A^k-D_k)=D_k+K_k\), proving
\cref{eq:DK-interval}.

Finally, subadditivity of \(\Psi\) gives
\[
 M^k+E_k=\Psi(C_k+D_k)
 \le\Psi(C_k)+\Psi(D_k)=M^k+\Psi(D_k),
\]
and \cref{eq:sharp-Psi} gives the last inequality.
\end{proof}

\begin{remark}[What the interval identities say]
Equation \cref{eq:E-interval} converts the archimedean convexity defect
\(D_k\) into an exact excess of 3-adic valuation weights on a consecutive
interval.  Equation \cref{eq:DK-interval} converts the return defect into a
corresponding 2-adic valuation sum.  A proof by this route would need more
than average-order information: it would have to exploit the special
endpoints \(\Phi(M^k)\), \(A^k\), \(M^k\), and \(\Psi(A^k)\).
\end{remark}

\subsection{Semigroup inequalities}

\begin{proposition}[Defect semigroup laws]\label{prop:defect-semigroup}
Set
\[
 d_k=\frac{D_k}{A^k},\qquad e_k=\frac{E_k}{M^k}.
\]
Then for all \(k,\ell\ge1\),
\begin{align}
 d_{k+\ell}&\le d_k+d_\ell-d_kd_\ell,
 \label{eq:d-semigroup}\\
 e_{k+\ell}&\le e_k+e_\ell+e_ke_\ell.
 \label{eq:e-semigroup}
\end{align}
Equivalently, \(1-d_k\) is supermultiplicative and \(1+e_k\) is
submultiplicative.
\end{proposition}

\begin{proof}
By \cref{eq:Phi-super},
\[
 \Phi(M^{k+\ell})\ge\Phi(M^k)\Phi(M^\ell).
\]
Divide by \(A^{k+\ell}\).  By \cref{eq:Psi-sub},
\[
 \Psi(A^{k+\ell})\le\Psi(A^k)\Psi(A^\ell),
\]
and division by \(M^{k+\ell}\) gives the second assertion.
\end{proof}

\section{The digit-difference polynomial and borrow rigidity}
\label{sec:digit-polynomial}

Let
\[
 M^k=\sum_{j=0}^{L_k}\epsilon_{k,j}2^j,
 \qquad
 A^k=\sum_{j=0}^{L_k}\delta_{k,j}3^j,
\]
where leading zero digits are inserted as needed.  The equality of the two
upper indices will be proved in \cref{sec:prefix}.  For the moment, take a
common sufficiently large upper index.

\begin{definition}[Difference polynomial]\label{def:Pk}
Define
\begin{equation}\label{eq:Pk}
 P_k(z)=\sum_j(\delta_{k,j}-\epsilon_{k,j})z^j
       \in\Z[z].
\end{equation}
\end{definition}

\begin{proposition}[Two evaluations of one word]\label{prop:Pk-eval}
For every \(k\ge1\),
\begin{equation}\label{eq:Pk-evals}
  P_k(3)=D_k,
  \qquad
  P_k(2)=E_k.
\end{equation}
If \(J_k=\deg P_k\), then its leading coefficient is positive.  More
precisely, it lies in \(\{1,2\}\).
\end{proposition}

\begin{proof}
The first identity is the definition of \(D_k\).  For the second,
\[
 P_k(2)=\Psi(A^k)-\Psi(\Phi(M^k))
       =\Psi(A^k)-M^k=E_k.
\]
If the leading coefficient at degree \(J\) were \(-1\), then even the
largest possible lower contribution would be
\[
 -3^J+2\sum_{j<J}3^j=-1<0,
\]
contradicting \(D_k>0\).  The coefficient cannot be less than \(-1\), and
its positive possibilities are \(1,2\).
\end{proof}

\begin{lemma}[Leading-coefficient lower bounds]\label{lem:lead-lower}
Let \(P(z)=\sum_{j=0}^Jc_jz^j\) with
\(c_j\in\{-1,0,1,2\}\) and positive leading coefficient.  Then
\begin{align}
 P(3)&\ge\frac{3^J+1}{2},\label{eq:P3-lower}\\
 P(2)&\ge1.\label{eq:P2-lower}
\end{align}
If \(c_J=2\), then
\begin{equation}\label{eq:P2-leading2}
 P(2)\ge2^J+1.
\end{equation}
\end{lemma}

\begin{proof}
The minimum at \(3\) occurs when all lower coefficients equal \(-1\):
\[
 3^J-\sum_{j<J}3^j=\frac{3^J+1}{2}.
\]
At \(2\), the same calculation gives \(1\) for leading coefficient \(1\)
and \(2^J+1\) for leading coefficient \(2\).
\end{proof}

The exceptional way in which \(P(2)\) can be small while \(P(3)\) is large
is completely rigid.

\begin{theorem}[Borrow-block lemma]\label{thm:borrow-block}
Let \(P(z)=\sum_{i=0}^Jc_iz^i\) be as in
\cref{lem:lead-lower}, and set \(E=P(2)>0\).

If \(c_J=1\), then
\begin{equation}\label{eq:borrow-expansion1}
 E=1+\sum_{i<J}(c_i+1)2^i.
\end{equation}
If \(c_J=2\), then
\begin{equation}\label{eq:borrow-expansion2}
 E=2^J+1+\sum_{i<J}(c_i+1)2^i.
\end{equation}
Consequently, if \(E<2^s\) with \(s\le J\), then necessarily
\(c_J=1\) and
\begin{equation}\label{eq:borrow-block-coeff}
 c_i=-1\qquad(s\le i<J).
\end{equation}
For the digit polynomial \(P_k\), this means
\begin{equation}\label{eq:borrow-block-digits}
 (\delta_{k,i},\epsilon_{k,i})=(0,1)
 \qquad(s\le i<J_k).
\end{equation}
\end{theorem}

\begin{proof}
Use \(\sum_{i<J}2^i=2^J-1\) to add and subtract the all-ones lower word.
Every coefficient \(c_i+1\) is nonnegative.  If the total is below \(2^s\),
all terms at positions \(i\ge s\) must vanish.  Since
\(c_i=\delta_{k,i}-\epsilon_{k,i}\), the value \(-1\) occurs only for the
pair \((0,1)\).
\end{proof}

\begin{corollary}[Congruence form of a long borrow block]\label{cor:borrow-cong}
Under the hypotheses of \cref{thm:borrow-block}, if \(E_k<2^s\) and
\(s\le J_k\), then there are integers \(u,v\) with
\begin{equation}\label{eq:uv-ranges}
 1\le u\le2^s,
 \qquad 0\le v<3^s,
\end{equation}
such that
\begin{equation}\label{eq:borrow-congruences}
 M^k\equiv-u\pmod{2^{J_k}},
 \qquad
 A^k\equiv v\pmod{3^{J_k}}.
\end{equation}
\end{corollary}

\begin{proof}
The binary bits in positions \(s,\ldots,J_k-1\) are all \(1\), so modulo
\(2^{J_k}\) the lower word is \(2^{J_k}-u\) with \(1\le u\le2^s\).
The ternary digits in the same positions vanish, leaving a residue below
\(3^s\).
\end{proof}

\subsection{The exact carry mask for the return defect}

\begin{theorem}[Carry-mask formula]\label{thm:carry-mask}
Let \(a=\sum_{j\ge0}\delta_j3^j\), and normalize the digit string
\(\delta_j\in\{0,1,2\}\) in base \(2\) by carry bits \(c_j\in\{0,1\}\):
\begin{equation}\label{eq:carry-rec}
 c_{-1}=0,
 \qquad
 \delta_j+c_{j-1}=\epsilon'_j+2c_j,
 \qquad \epsilon'_j\in\{0,1\}.
\end{equation}
Then
\begin{equation}\label{eq:K-carry-mask}
 \Phi(\Psi(a))-a=\sum_{j\ge0}c_j3^j.
\end{equation}
In particular, if the leading ternary digit of \(a\) at position \(L\) is
\(2\), then
\begin{equation}\label{eq:K-top2}
 \Phi(\Psi(a))-a\ge3^L.
\end{equation}
\end{theorem}

\begin{proof}
Equation \cref{eq:carry-rec} gives
\(\epsilon'_j-\delta_j=c_{j-1}-2c_j\).  Therefore
\[
 \sum_j(\epsilon'_j-\delta_j)3^j
 =3\sum_jc_j3^j-2\sum_jc_j3^j
 =\sum_jc_j3^j.
\]
If \(\delta_L=2\), then \(c_L=1\).
\end{proof}

This formula upgrades the inequality \(K_k\ge0\) to a finite two-state
transducer: \(K_k\) is literally the base-3 evaluation of its carry mask.

\section{Most-significant prefix separation}
\label{sec:prefix}

This section contains the central new structural theorem.

\subsection{Aligned leading positions}

Let
\begin{equation}\label{eq:L-u}
 L_k=\lfloor k\beta\rfloor,
 \qquad
 u_k=\{k\beta\}\in(0,1).
\end{equation}
The fractional part is never zero because \(\beta\notin\Q\).  Then
\begin{equation}\label{eq:mantissas}
 M^k=2^{L_k}2^{u_k},
 \qquad
 A^k=3^{L_k}3^{u_k}.
\end{equation}
Thus the most-significant binary and ternary digit positions are both
\(L_k\).  This exact alignment is special to the exponent pair
\((2^x,3^x)\).

For \(0<u<1\), put
\begin{equation}\label{eq:xy-mantissa}
 X(u)=2^u,
 \qquad Y(u)=3^u=X(u)^\vartheta.
\end{equation}
Write the fractional expansions
\begin{align}
 X(u)&=1+\sum_{t\ge1}b_t(u)2^{-t},
 &b_t(u)&\in\{0,1\},\label{eq:Xdigits}\\
 Y(u)&=h_0(u)+\sum_{t\ge1}a_t(u)3^{-t},
 &a_t(u)&\in\{0,1,2\},\label{eq:Ydigits}
\end{align}
using the nonterminating ambiguity only in the canonical terminating
fashion appropriate to the integer powers in \cref{eq:mantissas}.
Here \(h_0(u)=1\) for \(u<\log_3 2\) and \(h_0(u)=2\) otherwise.

\begin{lemma}[Cylinder-separation inequality]\label{lem:cylinder-sep}
For every integer \(t\ge1\),
\begin{equation}\label{eq:cylinder-sep}
  (1+2^{-t})^\vartheta>1+2\,3^{-t}.
\end{equation}
\end{lemma}

\begin{proof}
Bernoulli's inequality and \(2^{-t\vartheta}=3^{-t}\) give
\[
 (1+2^{-t})^\vartheta
 >1+\vartheta2^{-t}.
\]
It remains to note that
\[
 \vartheta2^{-t}>2\,3^{-t}
 \quad\Longleftrightarrow\quad
 \vartheta\left(\frac32\right)^t>2,
\]
and the left side is increasing in \(t\), while at \(t=1\) it is
\(3\vartheta/2>2\).
\end{proof}

\begin{theorem}[Prefix-separation theorem]\label{thm:prefix-separation}
Fix \(u\in(0,1)\) and compare the base-2 digit string of \(X(u)=2^u\) with
the base-3 digit string of \(Y(u)=3^u\), starting at their units digits and
moving right.

\begin{enumerate}[label=(\alph*)]
\item If \(u\ge\log_3 2\), the first digits are \(1\) and \(2\), so they
mismatch immediately.
\item If \(0<u<\log_3 2\), their common prefix is exactly
\begin{equation}\label{eq:common-prefix-form}
  1\,0\,0\cdots0.
\end{equation}
At the first mismatch after this prefix, the ternary digit is strictly
larger than the binary digit.
\end{enumerate}
Equivalently, after the common leading \(1\), the two strings can never
share a digit \(1\).
\end{theorem}

\begin{proof}
Part (a) is immediate.  For part (b), suppose the strings agree through the
first \(t-1\) fractional positions, all of which must initially be zeros.
If the binary digit at position \(t\) were \(1\), then
\(X(u)\ge1+2^{-t}\).  By \cref{lem:cylinder-sep},
\[
 Y(u)=X(u)^\vartheta>1+2\,3^{-t}.
\]
Given that the preceding ternary digits are zero, its digit at position
\(t\) must be \(2\), not \(1\).  If instead the ternary expansion becomes
nonzero first, then the binary digit is \(0\) and the ternary digit is
\(1\) or \(2\).  In every case the first mismatch has ternary digit larger
than binary digit.  Therefore no common \(1\) can occur after the leading
one, and all common intervening digits are zeros.
\end{proof}

\begin{corollary}[Exact first-mismatch depth]\label{cor:first-mismatch-depth}
For \(0<u<\log_3 2\), define
\begin{equation}\label{eq:t-of-u}
 t(u)=\min\{t\ge1:3^u-1\ge3^{-t}\}
     =\left\lceil\log_3\frac1{3^u-1}\right\rceil.
\end{equation}
Then the common prefix has the leading \(1\) followed by exactly
\(t(u)-1\) zeros, and the first mismatch is at fractional position \(t(u)\).
For the integer powers \(M^k,A^k\), the degree of \(P_k\) is
\begin{equation}\label{eq:Jk-formula}
 J_k=
 \begin{cases}
 L_k,&u_k\ge\log_3 2,\\
 L_k-t(u_k),&0<u_k<\log_3 2.
 \end{cases}
\end{equation}
The leading coefficient of \(P_k\) is exactly the positive digit difference
at this first mismatch.
\end{corollary}

\begin{proof}
The first nonzero ternary fractional digit occurs exactly when the ternary
remainder \(3^u-1\) reaches \(3^{-t}\).  The rest follows from
\cref{thm:prefix-separation} after scaling by \(2^{L_k}\) and \(3^{L_k}\).
\end{proof}

\begin{remark}[A symbolic rigidity not visible to real approximation]
The theorem is stronger than the convexity inequality
\(\Phi(M^k)<A^k\).  It identifies the most-significant mismatch and proves
that every aligned digit above it is forced: a single \(1\), then only
zeros.  No arbitrary common word is possible.  Thus the cross-base digit
comparison has a one-parameter cylinder geometry rather than a full
shift-space geometry.
\end{remark}

\subsection{Separation near the top}

Without any transcendence estimate, equidistribution already gives a
statistical form of top separation.  An effective all-\(k\) form follows
from linear forms in logarithms.

\begin{proposition}[Density-one bounded prefix depth]\label{prop:prefix-density}
For every \(\eta\in(0,1)\), there is an integer \(T_\eta\) such that the set
of \(k\) for which
\begin{equation}\label{eq:J-density}
 J_k\ge L_k-T_\eta
\end{equation}
has natural density at least \(1-\eta\).
\end{proposition}

\begin{proof}
The irrationality of \(\beta\) implies that \(u_k=\{k\beta\}\) is uniformly
distributed modulo one.  Outside an interval \((0,\delta)\) of measure
\(\delta\), the quantity \(3^{u_k}-1\) is bounded below, so
\cref{eq:t-of-u} is bounded above by a constant depending only on
\(\delta\).  Choose \(\delta=\eta\).
\end{proof}

\begin{theorem}[Effective logarithmic prefix bound]\label{thm:prefix-log}
For a fixed hypothetical pair \((M,A)\), there is an effectively computable
constant \(C_M>0\) such that, for every \(k\ge2\),
\begin{equation}\label{eq:t-log-bound}
 t(u_k)\le C_M\log k
 \quad\text{whenever }u_k<\log_3 2.
\end{equation}
Consequently,
\begin{equation}\label{eq:J-log-bound}
 J_k\ge L_k-C_M\log k.
\end{equation}
\end{theorem}

\begin{proof}
The nonzero linear form
\begin{equation}\label{eq:linear-form-prefix}
 \Lambda_k=k\log M-L_k\log2=u_k\log2
\end{equation}
has algebraic arguments \(M\) and \(2\).  An explicit lower bound for a
linear form in two logarithms (for example, a specialization of Matveev's
theorem) yields
\begin{equation}\label{eq:u-baker}
 u_k\ge k^{-C_0(M)}
\end{equation}
for \(k\ge2\), after enlarging the effective constant.  On
\(0<u<\log_3 2\), the mean value theorem gives
\(3^u-1\ge(\log3)u\).  Substitute this into \cref{eq:t-of-u}:
\[
 t(u_k)
 \le1+\log_3\frac1{(\log3)u_k}
 \le C_M\log k.
\]
Equation \cref{eq:J-log-bound} follows from \cref{eq:Jk-formula}.
\end{proof}

\begin{corollary}[Exponential size of the direct defect]\label{cor:D-from-J}
For a fixed hypothetical pair there is an effective \(C>0\) such that
\begin{equation}\label{eq:D-J-bound}
 D_k\ge\frac{3^{J_k}+1}{2}
      \gg_{M,A}\frac{A^k}{k^C}
 \qquad(k\ge2).
\end{equation}
\end{corollary}

\begin{proof}
Apply \cref{lem:lead-lower} and
\(3^{J_k}\ge3^{L_k}k^{-C'}\).  Since
\(A^k=3^{L_k+u_k}<3^{L_k+1}\), the result follows.
\end{proof}

\section{Macroscopic defect laws}
\label{sec:macro}

The preceding digit theorem identifies the mismatch position.  Independent
convexity estimates produce explicit lower envelopes for the relative
defects.

\subsection{Explicit envelope functions}

Define, for \(0<u<1\),
\begin{equation}\label{eq:d-envelope}
 \mathfrak d(u)
 =\frac{3^u-1-(2^u-1)^\vartheta}{3^u}.
\end{equation}
For \(y=3^u\) and \(h=\lfloor y\rfloor\in\{1,2\}\), define
\begin{equation}\label{eq:e-envelope}
 \mathfrak e(u)
 =\frac{h+(y-h)^\varrho-2^u}{2^u}.
\end{equation}
Both functions are positive on \((0,1)\).  They have positive left limits
at \(1\), and their only vanishing endpoint is \(u=0\).

\begin{theorem}[Relative defect envelopes]\label{thm:defect-envelopes}
For every \(k\ge1\),
\begin{align}
 \frac{D_k}{A^k}&\ge\mathfrak d(u_k),
 \label{eq:D-envelope}\\
 \frac{E_k}{M^k}&\ge\mathfrak e(u_k).
 \label{eq:E-envelope}
\end{align}
\end{theorem}

\begin{proof}
Write \(M^k=2^{L_k}+r\), where
\(r/2^{L_k}=2^{u_k}-1\).  By
\(\Phi(r)\le r^\vartheta\),
\[
 \Phi(M^k)
 =3^{L_k}+\Phi(r)
 \le3^{L_k}\bigl(1+(2^{u_k}-1)^\vartheta\bigr).
\]
Subtract from \(A^k=3^{L_k}3^{u_k}\) and divide by \(A^k\).

For the second estimate, write
\(A^k=h3^{L_k}+s\), where
\(s/3^{L_k}=3^{u_k}-h\).  By
\(\Psi(s)\ge s^\varrho\),
\[
 \Psi(A^k)
 =h2^{L_k}+\Psi(s)
 \ge2^{L_k}\bigl(h+(3^{u_k}-h)^\varrho\bigr).
\]
Subtract \(M^k=2^{L_k}2^{u_k}\) and divide by it.
Strict positivity is the strict convexity or concavity already used in
\cref{prop:binary-defect,prop:ternary-defect}.
\end{proof}

\begin{corollary}[Positive-density macroscopic defects]\label{cor:density-macro}
For every \(\eta\in(0,1)\), there are constants
\(c_D(\eta),c_E(\eta)>0\) such that each of the inequalities
\begin{equation}\label{eq:density-macro}
 D_k\ge c_D(\eta)A^k,
 \qquad
 E_k\ge c_E(\eta)M^k
\end{equation}
holds simultaneously on a set of \(k\) of natural density at least
\(1-\eta\).
\end{corollary}

\begin{proof}
On the compact interval \([\eta,1]\), after taking left limits at \(1\),
both envelope functions have a positive minimum.  Uniform distribution of
\(u_k\) gives the asserted density.
\end{proof}

\begin{corollary}[Positive Ces\`aro means]\label{cor:Cesaro}
One has
\begin{align}
 \liminf_{N\to\infty}\frac1N\sum_{k\le N}\frac{D_k}{A^k}
 &\ge\int_0^1\mathfrak d(u)\,du>0,
 \label{eq:Cesaro-D}\\
 \liminf_{N\to\infty}\frac1N\sum_{k\le N}\frac{E_k}{M^k}
 &\ge\int_0^1\mathfrak e(u)\,du>0.
 \label{eq:Cesaro-E}
\end{align}
\end{corollary}

\begin{proof}
The envelope functions are bounded and Riemann integrable, including at the
single branch point \(3^u=2\).  Apply Weyl equidistribution to
\(u_k=\{k\beta\}\) and use \cref{thm:defect-envelopes}.
\end{proof}

\subsection{Effective all-exponent lower bounds}

\begin{theorem}[Polynomial-loss macroscopic lower bounds]\label{thm:all-k-macro}
For a fixed hypothetical pair \((M,A)\), there are effectively computable
constants \(C_1,C_2>0\) and \(c_1,c_2>0\) such that, for all \(k\ge2\),
\begin{equation}\label{eq:all-k-macro}
 D_k\ge c_1\frac{A^k}{k^{C_1}},
 \qquad
 E_k\ge c_2\frac{M^k}{k^{C_2}}.
\end{equation}
\end{theorem}

\begin{proof}
By \cref{eq:u-baker}, \(u_k\ge k^{-C_0}\).  Near zero,
\cref{prop:binary-defect} applied to the mantissa gives
\(\mathfrak d(u)\gg u\).  For \(\mathfrak e\), when \(u\) is small the
leading ternary digit is \(1\), and \cref{eq:ternary-defect-d1} gives
\(\mathfrak e(u)\gg(3^u-1)^\varrho\gg u^\varrho\).  Away from zero, both
functions have a positive lower bound.  Insert the lower bound for \(u_k\)
into \cref{eq:D-envelope,eq:E-envelope}.
\end{proof}

\begin{corollary}[Only logarithmic catastrophic borrowing]\label{cor:borrow-log}
Let \(s_k=\lfloor\log_2E_k\rfloor+1\).  If the borrow block in
\cref{thm:borrow-block} is nonempty, then its length satisfies
\begin{equation}\label{eq:borrow-log}
 J_k-s_k=O_{M,A}(\log k).
\end{equation}
\end{corollary}

\begin{proof}
By \cref{thm:all-k-macro},
\(s_k\ge\log_2M^k-O(\log k)=L_k-O(\log k)\).  Since \(J_k\le L_k\), the
claim follows.
\end{proof}

\begin{remark}[A useful negative conclusion]
The defects are not candidates for bounded-gap arguments.  They are
exponentially large for almost all exponents and, with only a polynomial
loss, for every exponent.  Any successful use of \(D_k,E_k\) must exploit
factorization, congruence, or digit structure rather than small absolute
size.
\end{remark}

\section{Finite-place residue transfer}
\label{sec:residues}

\begin{proposition}[Residue transfer]\label{prop:residue-transfer}
For all \(r\ge1\),
\begin{align}
 n\equiv m\pmod{2^r}
 &\Longrightarrow
 \Phi(n)\equiv\Phi(m)\pmod{3^r},
 \label{eq:Phi-residue}\\
 a\equiv b\pmod{3^r}
 &\Longrightarrow
 \Psi(a)\equiv\Psi(b)\pmod{2^r}.
 \label{eq:Psi-residue}
\end{align}
\end{proposition}

\begin{proof}
Only the first \(r\) source digits contribute modulo the \(r\)-th power of
the target prime.
\end{proof}

The primitive normalization imported from the supplied report implies that
a hypothetical pair falls into three local branches: both \(M\) and \(A\)
are units at \(2\) and \(3\), respectively; \(M\) is odd while \(3\mid A\);
or \(2\mid M\) while \(3\nmid A\).  The branch \(2\mid M\), \(3\mid A\)
is excluded by that normalization.

\begin{proposition}[Three recurrence branches]\label{prop:three-branches}
Fix \(r\ge1\).
\begin{enumerate}[label=(\roman*)]
\item If \(M\) is odd and \(3\nmid A\), there are infinitely many \(k\)
for which
\begin{equation}\label{eq:branch-unit}
 D_k\equiv0\pmod{3^r},
 \qquad E_k\equiv0\pmod{2^r}.
\end{equation}
\item If \(M\) is odd and \(3\mid A\), there are infinitely many \(k\)
for which
\begin{equation}\label{eq:branch-3}
 D_k\equiv-1\pmod{3^r},
 \qquad E_k\equiv-1\pmod{2^r}.
\end{equation}
\item If \(2\mid M\) and \(3\nmid A\), there are infinitely many \(k\)
for which
\begin{equation}\label{eq:branch-2}
 D_k\equiv1\pmod{3^r},
 \qquad E_k\equiv1\pmod{2^r}.
\end{equation}
\end{enumerate}
\end{proposition}

\begin{proof}
In (i), choose \(k\) divisible by the multiplicative orders of \(M\) modulo
\(2^r\) and of \(A\) modulo \(3^r\).  Then
\(M^k\equiv1\pmod{2^r}\), \(A^k\equiv1\pmod{3^r}\), and use
\cref{prop:residue-transfer}.  In (ii), impose
\(M^k\equiv1\pmod{2^r}\) and take \(k\) large enough that
\(A^k\equiv0\pmod{3^r}\).  Case (iii) is symmetric.
\end{proof}

\begin{remark}[The scale obstruction]
For a common exponent \(k\), lifting-the-exponent estimates on the unit side
usually provide only \(r=O(\log k)\).  The other side may have linearly
growing valuation, but the shared modulus is limited by the unit condition.
Thus the forced product \(2^r3^r\) is only polynomial in \(k\), whereas
\cref{thm:all-k-macro} places the defects on an exponential scale.  A direct
product-formula contradiction misses an exponential factor.
\end{remark}

\section{Pure-radical norm gaps}
\label{sec:radical-norm}

We now use the imported normalization \cref{eq:radical-import}.  Put
\begin{equation}\label{eq:alpha-radical}
 \alpha=\log_2w=\log_3\eta.
\end{equation}
For integers \(q\ge1\), \(p\in\Z\), define
\begin{equation}\label{eq:gDm}
 g=\gcd(d,q),
 \qquad D=\frac d g,
 \qquad m=\frac q g.
\end{equation}

\begin{proposition}[Degree of a powered radical]\label{prop:powered-radical}
Under the exact-degree hypothesis in \cref{eq:radical-import},
\(\eta^q\) has degree \(D=d/g\), and its minimal polynomial is
\begin{equation}\label{eq:minpoly-powered}
 X^D-\left(\frac{A}{3^a}\right)^m.
\end{equation}
\end{proposition}

\begin{proof}
Choose integers \(r,s\) such that \(rq+sd=g\).  Then
\[
 \eta^g=(\eta^q)^r(\eta^d)^s\in\Q(\eta^q),
\]

while \(\eta^q=(\eta^g)^{q/g}\), so
\(\Q(\eta^q)=\Q(\eta^g)\).  The exact pure-radical degree, equivalently the
standard irreducibility criterion for \(X^d-A/3^a\), passes to the divisor
\(D\); hence \([\Q(\eta^g):\Q]=D\).  The displayed polynomial has degree
\(D\) and vanishes at \(\eta^q\), so it is minimal.
\end{proof}

\begin{corollary}[Integral norm numerator]\label{cor:norm-numerator}
For every integer \(q\ge1\) and every integer \(p\) with \(am+pD\ge0\), the
cleared norm of \(\eta^q-3^p\) is the nonzero integer
\begin{equation}\label{eq:Rqp}
 R(q,p)=\left|A^m-3^{am+pD}\right|
 =3^{am+pD}\left|3^{D(q\alpha-p)}-1\right|.
\end{equation}
\end{corollary}

\begin{proof}
Take the norm using \cref{eq:minpoly-powered} and multiply by \(3^{am}\).
The second equality follows from
\(A^m=3^{am+Dq\alpha}\).
\end{proof}

Let \(p_n/q_n\) be consecutive continued-fraction convergents to \(\alpha\),
and define \(g_n,D_n,m_n\) by \cref{eq:gDm}.  The exponent vector in
\cref{eq:Rqp} is
\begin{equation}\label{eq:radical-vector}
 \mathbf v_n=
 \left(m_n,\,am_n+p_nD_n\right).
\end{equation}

\begin{proposition}[Bounded determinant radical ladder]\label{prop:radical-det}
For adjacent convergents,
\begin{equation}\label{eq:radical-det}
 \left|\det(\mathbf v_n,\mathbf v_{n+1})\right|
 =\frac{d}{g_ng_{n+1}},
\end{equation}
which is a positive divisor of \(d\).
\end{proposition}

\begin{proof}
The determinant equals
\[
 \frac{d}{g_ng_{n+1}}
 |q_np_{n+1}-q_{n+1}p_n|.
\]
The final factor is \(1\).  Since adjacent denominators are coprime,
\(g_n\) and \(g_{n+1}\) are coprime divisors of \(d\), so their product
divides \(d\).
\end{proof}

The determinant-gcd lemma imported from the supplied report therefore gives
a fixed finite bound for adjacent gcds after removing primes dividing
\(3A\).  In the branch \(a>0\), primitive normalization gives \(3\nmid A\),
and the full adjacent gcd is bounded by a fixed integer.  In every branch,
an S-unit finiteness theorem implies that the union of prime divisors of
\(R(q_n,p_n)\) cannot remain finite.  This proves unbounded prime support
along the ladder, but not a new prime at every step.

\subsection{A rational secant to \(\vartheta\)}

Set \(r_n=am_n+p_nD_n\) and
\(B_n=|w^{q_n}-2^{p_n}|\).  Whenever \(B_n\ne0\), define
\begin{equation}\label{eq:Qn-secant}
 Q_n=\frac{2^{p_n}R(q_n,p_n)}{D_n3^{r_n}B_n}.
\end{equation}
If \(\epsilon_n=q_n\alpha-p_n\), cancellation gives the exact identity
\begin{equation}\label{eq:Qn-cancel}
 Q_n=\frac{|3^{D_n\epsilon_n}-1|}
           {D_n|2^{\epsilon_n}-1|}
 \longrightarrow\vartheta.
\end{equation}
This produces rational approximants to \(\vartheta\) from two coupled
integer gaps.  Their heights, however, are exponential in \(q_n\), while
the error is only of continued-fraction order.  Thus the construction does
not contradict known irrationality measures; it is recorded because the
same cancellation becomes more structural in the next section.

\section{The synchronized determinant-one gap ladder}
\label{sec:gap-ladder}

The previous ladder has determinant dividing \(d\).  A better choice of the
approximated logarithm makes the determinant exactly one and synchronizes
the two bases with the \emph{same} logarithmic error.

Set
\begin{equation}\label{eq:gamma}
 \gamma=d\log_2w.
\end{equation}
By \cref{eq:radical-import},
\begin{equation}\label{eq:beta-gamma}
 \beta=\log_2M=\log_3A=a+\gamma.
\end{equation}
Let \(p_n/q_n\) be the ordinary continued-fraction convergents to
\(\gamma\), and put
\begin{equation}\label{eq:eps-r}
 \epsilon_n=q_n\gamma-p_n,
 \qquad r_n=aq_n+p_n.
\end{equation}
Define the simultaneous gaps
\begin{equation}\label{eq:BR-def}
 B_n=\left|w^{dq_n}-2^{p_n}\right|,
 \qquad
 R_n=\left|A^{q_n}-3^{r_n}\right|.
\end{equation}

\begin{theorem}[Same-error synchronization]\label{thm:same-error}
The gaps satisfy
\begin{align}
 B_n&=2^{p_n}\left|2^{\epsilon_n}-1\right|,
 \label{eq:B-same-error}\\
 R_n&=3^{r_n}\left|3^{\epsilon_n}-1\right|.
 \label{eq:R-same-error}
\end{align}
The adjacent exponent vectors \((q_n,r_n)\) have determinant one:
\begin{equation}\label{eq:det-one}
 |q_nr_{n+1}-q_{n+1}r_n|=1.
\end{equation}
\end{theorem}

\begin{proof}
The first pair of identities follows from
\(w^{dq_n}=2^{q_n\gamma}=2^{p_n+\epsilon_n}\) and
\(A^{q_n}=3^{q_n(a+\gamma)}=3^{r_n+\epsilon_n}\).  For the determinant,
\[
 q_nr_{n+1}-q_{n+1}r_n
 =q_np_{n+1}-q_{n+1}p_n,
\]
because the terms containing \(a\) cancel.  Adjacent continued-fraction
convergents have determinant \(\pm1\).
\end{proof}

\begin{corollary}[Sharp adjacent coprimality]\label{cor:adjacent-coprime}
Let
\[
 G_n=M^{q_n}-2^{r_n}
     =\pm2^{aq_n}B_n.
\]
The odd parts of \(G_n\) and \(G_{n+1}\) are coprime.  The parts of
\(R_n,R_{n+1}\) away from \(3\) have gcd at most \(2\).  If \(3\nmid A\),
then
\begin{equation}\label{eq:R-gcd2}
 \gcd(R_n,R_{n+1})\le2.
\end{equation}
\end{corollary}

\begin{proof}
Apply the determinant-gcd lemma for near-power gaps with determinant one.
For bases \((M,2)\), the prime-to-\(2M\) gcd divides both \(M-1\) and
\(2-1=1\); primes dividing the odd part of \(M\) cannot divide the gap, so
the odd parts are coprime.  For bases \((A,3)\), the prime-to-\(3A\) gcd
divides \(\gcd(A-1,3-1)\le2\); primes dividing \(A\) but not \(3\) cannot
divide \(A^{q_n}-3^{r_n}\).  If \(3\nmid A\), no 3-adic factor remains.
\end{proof}

\subsection{The exact gap-descent square}

The errors \(\epsilon_n\) alternate in sign.  Restrict to the infinite
subsequence \(\epsilon_n>0\).  For all sufficiently large indices on this
subsequence,
\begin{equation}\label{eq:gap-small}
 0<B_n<2^{p_n},
 \qquad
 0<R_n<3^{r_n}.
\end{equation}
Then
\begin{align}
 M^{q_n}&=2^{r_n}+2^{aq_n}B_n,
 \label{eq:M-gap-decomp}\\
 A^{q_n}&=3^{r_n}+R_n.
 \label{eq:A-gap-decomp}
\end{align}
The two summands in each line occupy disjoint digit blocks.

\begin{theorem}[Gap-descent square]\label{thm:gap-descent}
For all sufficiently large \(n\) with \(\epsilon_n>0\),
\begin{align}
 \Phi(M^{q_n})
 &=3^{r_n}+3^{aq_n}\Phi(B_n),
 \label{eq:Phi-gap-decomp}\\
 \Psi(A^{q_n})
 &=2^{r_n}+\Psi(R_n),
 \label{eq:Psi-gap-decomp}
\end{align}
and therefore
\begin{align}
 D_{q_n}
 &=R_n-3^{aq_n}\Phi(B_n)>0,
 \label{eq:D-gap-descent}\\
 E_{q_n}
 &=\Psi(R_n)-2^{aq_n}B_n>0.
 \label{eq:E-gap-descent}
\end{align}
\end{theorem}

\begin{proof}
Because \(B_n<2^{p_n}\), the shifted integer \(2^{aq_n}B_n\) uses binary
positions strictly below \(r_n=aq_n+p_n\).  Hence no carry occurs in
\cref{eq:M-gap-decomp}, and digit transfer gives
\cref{eq:Phi-gap-decomp}.  Similarly, \(R_n<3^{r_n}\) gives
\cref{eq:Psi-gap-decomp}.  Subtract from \cref{eq:A-gap-decomp} and
\cref{eq:M-gap-decomp}; positivity is \cref{prop:defect-positive}.
\end{proof}

The theorem is an exact commutative-square phenomenon: the same convergent
error produces gaps \(B_n,R_n\); the transfer maps then compare the
large-power defects to transferred versions of the small gaps.

\subsection{Why the square is not contractive}

\begin{theorem}[Scale-mismatch no-go theorem]\label{thm:scale-mismatch}
Along the positive-error subsequence,
\begin{align}
 \frac{3^{aq_n}\Phi(B_n)}{R_n}
 &=O_{M,A}(\epsilon_n^{\vartheta-1})\longrightarrow0,
 \label{eq:ratio-one}\\
 \frac{2^{aq_n}B_n}{\Psi(R_n)}
 &=O_{M,A}(\epsilon_n^{1-\varrho})\longrightarrow0.
 \label{eq:ratio-two}
\end{align}
Consequently,
\begin{equation}\label{eq:defect-asymp-gaps}
 D_{q_n}\sim R_n,
 \qquad
 E_{q_n}\sim\Psi(R_n).
\end{equation}
Both asymptotics are in the ratio sense: each quotient of the left side by
the right side tends to \(1\).
\end{theorem}

\begin{proof}
For small positive \(\epsilon\),
\(2^\epsilon-1\asymp\epsilon\) and
\(3^\epsilon-1\asymp\epsilon\).  Thus
\[
 B_n\asymp2^{p_n}\epsilon_n,
 \qquad
 R_n\asymp3^{r_n}\epsilon_n.
\]
By \(\Phi(B_n)\le B_n^\vartheta\),
\[
 3^{aq_n}\Phi(B_n)
 \ll3^{aq_n}(2^{p_n}\epsilon_n)^\vartheta
 =3^{r_n}\epsilon_n^\vartheta,
\]
which proves \cref{eq:ratio-one}.  By
\(\Psi(R_n)\ge R_n^\varrho\),
\[
 \Psi(R_n)\gg(3^{r_n}\epsilon_n)^\varrho
 =2^{r_n}\epsilon_n^\varrho,
\]
whereas \(2^{aq_n}B_n\asymp2^{r_n}\epsilon_n\).  This proves
\cref{eq:ratio-two}.  Insert the limits into
\cref{eq:D-gap-descent,eq:E-gap-descent}.
\end{proof}

\begin{warning}
The phrase ``gap descent'' does not mean a descent in size.  The exact
identities are structurally strong, but the transferred correction is
asymptotically negligible relative to the target gap.  Any proof that
silently assumes comparable scales reverses the actual asymptotics.
\end{warning}

\section{What the new framework rules out}
\label{sec:nogo}

The results above close several tempting but invalid lines of attack.

\begin{enumerate}
\item \textbf{Bounded or subexponential defects.}
By \cref{thm:all-k-macro}, both \(D_k\) and \(E_k\) are exponential up to a
polynomial factor.  A contradiction cannot come from showing merely that
they are nonzero integers.

\item \textbf{Long common arbitrary prefixes.}
By \cref{thm:prefix-separation}, the only common most-significant word is
\(1\) followed by zeros.  Any heuristic based on random agreement of
binary and ternary words is modeling the wrong symbolic system.

\item \textbf{Linear-length top agreement.}
By \cref{thm:prefix-log}, the common zero run has length only
\(O(\log k)\) for a fixed hypothetical pair.

\item \textbf{Uncontrolled catastrophic borrow.}
A very small value of \(P_k(2)\) forces a unique block of digit pairs
\((0,1)\), and \cref{cor:borrow-log} limits such a block to logarithmic
length once the all-\(k\) defect lower bound is used.

\item \textbf{A naive local product formula.}
The strongest immediately synchronized residue powers grow only
polynomially in \(k\), while the defects are exponentially large.

\item \textbf{A contracting continued-fraction descent.}
The exact square in \cref{thm:gap-descent} has the opposite scale ordering;
\cref{thm:scale-mismatch} quantifies the failure.

\item \textbf{A finite global prime set for the radical gaps.}
S-unit finiteness forces unbounded prime support, though it does not by
itself force a contradiction with the hypothetical pair.
\end{enumerate}

These eliminations are useful progress: they replace several qualitative
hopes by exact asymptotic obstructions and identify where genuinely new
arithmetic information must enter.

\section{Two precise closing targets}
\label{sec:targets}

The framework suggests two non-equivalent routes to a proof.  Each is
formulated so that it can be attacked independently and tested
computationally.

\subsection{Target A: prime transfer across the gap-descent square}

The determinant-one ladder gives adjacent coprimality, while
\cref{thm:gap-descent} links the gap \(R_n\) to \(\Phi(B_n)\) and
\(\Psi(R_n)\) to \(B_n\).  What is missing is a mechanism that forces a
substantial prime factor to survive digit transfer.

\begin{problem}[Prime-transfer problem]\label{prob:prime-transfer}
Let \(B,R\) be synchronized by
\[
 B=2^p(2^\epsilon-1),
 \qquad
 R=3^r(3^\epsilon-1),
\]
with integral \(B,R\), \(0<\epsilon\ll1\), and with the exponent vectors
belonging to a determinant-one continued-fraction ladder.  Prove that one of
the following must occur infinitely often:
\begin{enumerate}[label=(\alph*)]
\item a prime outside a fixed finite set divides both \(R\) and a controlled
arithmetic function of \(\Phi(B)\);
\item a prime outside a fixed finite set divides both \(B\) and a controlled
arithmetic function of \(\Psi(R)\);
\item the carry masks of \(\Phi(B)\) or \(\Psi(R)\) have complexity large
enough to contradict the two-term identities
\cref{eq:D-gap-descent,eq:E-gap-descent}.
\end{enumerate}
\end{problem}

A sufficiently uniform result of this kind would clash with adjacent
coprimality in \cref{cor:adjacent-coprime}.  Merely proving that the prime
support of \(R_n\) is unbounded is not enough; the transfer must interact
with neighboring indices or with the fixed bases.

\subsection{Target B: endpoint valuation excess on special intervals}

Equation \cref{eq:E-interval} may be rewritten as
\begin{equation}\label{eq:valuation-excess}
 \sum_{C_k<n\le A^k}\left(2^{v_3(n)}-1\right)=D_k-E_k.
\end{equation}
The interval length \(D_k\) and endpoint \(C_k=\Phi(M^k)\) are not generic.
Likewise, \cref{eq:DK-interval} is a special 2-adic valuation sum.

\begin{problem}[Endpoint-rigidity problem]\label{prob:endpoint-rigidity}
Prove an endpoint-sensitive estimate for one of the sums
\begin{align*}
 \sum_{\Phi(M^k)<n\le A^k}2^{v_3(n)},
 \qquad
 \sum_{M^k<n\le\Psi(A^k)}3^{v_2(n)},
\end{align*}
that uses the prefix-separation pattern of \cref{thm:prefix-separation} and
is stronger by an exponential factor than estimates based only on interval
length.
\end{problem}

The need for an exponential improvement is forced by the scale comparison in
\cref{sec:residues}.  Average-order valuation estimates alone cannot close
this target.

\subsection{A falsifiable symbolic criterion}

\begin{criterion}[Carry-word criterion]\label{crit:carry-word}
The conjecture would follow if one could prove that, for every pair of
integers \(M,A>1\) with \(\log_2M=\log_3A\notin\Q\), the digit-difference
polynomials \(P_k\) cannot simultaneously satisfy all three properties for
all \(k\):
\begin{enumerate}[label=(\roman*)]
\item their common most-significant prefix is \(1\,0^{O(\log k)}\);
\item their leading mismatch is positive and every near-cancellation at
\(z=2\) is carried by a \((0,1)\)-block of length \(O(\log k)\);
\item the two exact valuation interval identities
\cref{eq:E-interval,eq:DK-interval} hold with endpoints generated by the
same multiplicative orbit.
\end{enumerate}
\end{criterion}

The criterion deliberately packages only proved necessary properties.  It
is suitable for a finite-state or S-adic attack because every unbounded
feature has already been reduced to logarithmic depth.

\section{Computational reconnaissance program}
\label{sec:computation}

A numerical search cannot test a genuine counterexample without knowing one,
but it can test candidate closing lemmas on near-miss pairs and on exact
synchronized gaps.  The following experiments are mathematically
well-posed and avoid floating-point ambiguity by using integer arithmetic.

\subsection{Exact map and carry tests}

For integers \(M\) in a prescribed range, choose
\(A=\lfloor M^\vartheta\rceil\) only as a near-miss model.  For
\(1\le k\le K\), compute:
\begin{enumerate}
\item \(\Phi(M^k),\Psi(A^k),D_k,E_k,K_k\);
\item the degree \(J_k\), leading mismatch type, and borrow depth;
\item the carry mask \(c_j\) from \cref{eq:carry-rec};
\item the valuation sums in \cref{eq:E-interval,eq:DK-interval};
\item prime supports and adjacent gcds.
\end{enumerate}
For a near-miss pair the exact positivity relations need not all hold; this
is useful, because it reveals which patterns depend on exact exponent
matching and which are generic digit phenomena.

\subsection{Synchronized convergent tests}

For a chosen primitive tuple \((a,d,w)\), compute convergents to
\(\gamma=d\log_2w\) with certified interval arithmetic.  Retain those for
which \(B_n,R_n\) in \cref{eq:BR-def} are exact integers with the required
sign.  Then measure:
\begin{equation}\label{eq:computational-ratios}
 \frac{3^{aq_n}\Phi(B_n)}{R_n},
 \qquad
 \frac{2^{aq_n}B_n}{\Psi(R_n)},
 \qquad
 \gcd(R_n,R_{n+1}),
 \qquad
 \gcd(B_n,B_{n+1})_{\mathrm{odd}}.
\end{equation}
The predicted power laws in \cref{thm:scale-mismatch} provide direct
regression checks.  More importantly, one can search for stable residue or
prime-sharing phenomena not explained by determinant one.

\subsection{Search for a transfer invariant}

For each prime \(\ell\notin\{2,3\}\), tabulate the finite maps
\[
 n\pmod{2^r}\mapsto\Phi(n)\pmod\ell,
 \qquad
 a\pmod{3^r}\mapsto\Psi(a)\pmod\ell
\]
for the smallest \(r\) at which they stabilize on the relevant power orbit.
The objective is not to assume periodicity in the wrong modulus, but to
identify a finite quotient on which multiplication by \(M\) and \(A\)
interacts with the digit transducers.  Any invariant that couples two
adjacent convergents would be a candidate input for
\cref{prob:prime-transfer}.

\section{Lean formalization blueprint}
\label{sec:lean}

The new framework is unusually well suited to formalization because its
core is finite-digit algebra.  A practical dependency graph is as follows.

\subsection{Module structure}

\begin{longtable}{p{0.29\textwidth}p{0.64\textwidth}}
\toprule
Module & Principal content\\
\midrule
\endhead
\texttt{DigitTransfer.Basic} & Binary and ternary digit lists; definitions
of \(\Phi,\Psi\); radix recurrences; finite support.\\
\texttt{DigitTransfer.Carry} & Carry normalization; super/subadditivity;
composition; carry-mask theorem.\\
\texttt{DigitTransfer.Valuation} & Increment identities and summatory
formulas using \(v_2,v_3\).\\
\texttt{DigitTransfer.Bounds} & Convexity/concavity, sharp constants
\(1/2,2\), equality cases.\\
\texttt{AlaogluErdos.Defects} & Definitions of \(C_k,D_k,E_k,K_k\),
positivity, interval identities, semigroup inequalities.\\
\texttt{AlaogluErdos.Words} & Difference polynomial, leading sign,
borrow-block lemma, congruence corollary.\\
\texttt{AlaogluErdos.Prefix} & Mantissa alignment, cylinder-separation
inequality, prefix-separation theorem, first-mismatch formula.\\
\texttt{AlaogluErdos.Density} & Irrational rotation and positive-density
macroscopic defect bounds.\\
\texttt{AlaogluErdos.Congruences} & Residue transfer and the three local
branches.\\
\texttt{AlaogluErdos.Gaps} & Continued-fraction convergents, determinant
one, exact gap-descent identities.\\
\bottomrule
\end{longtable}

\subsection{Suggested theorem statements}

The following pseudo-Lean signatures indicate the intended granularity.
They are not claimed to compile verbatim.

\begin{verbatim}
def phi (n : Nat) : Nat :=
  (Nat.digits 2 n).eval 3

def psi (n : Nat) : Nat :=
  (Nat.digits 3 n).eval 2

theorem psi_phi (n : Nat) : psi (phi n) = n

theorem phi_sub_phi_pred (n : Nat) (hn : 0 < n) :
  phi n - phi (n - 1) = (3 ^ padicValNat 2 n + 1) / 2

theorem defect_interval (k : Nat) (hk : 0 < k) :
  E M A k = 2 * D M A k -
    \sum n in Ioc (phi (M^k)) (A^k), 2 ^ padicValNat 3 n

theorem prefix_separation (u : Real) (hu0 : 0 < u) (hu1 : u < 1) :
  CommonPrefix (base2Mantissa u) (base3Mantissa u) =
    [1] ++ List.replicate (prefixZeroLength u) 0
\end{verbatim}

\subsection{Formalization order}

The recommended order is:
\[
 \texttt{Basic}
 \to\texttt{Carry}
 \to\texttt{Bounds}
 \to\texttt{Defects}
 \to\texttt{Words}
 \to\texttt{Prefix}.
\]
This proves the entirely elementary core before importing analytic number
theory.  The Matveev-based \(O(\log k)\) theorem should initially be
formalized as a theorem parameterized by a lower-bound hypothesis for the
linear form.  A later module can connect that hypothesis to an available
formal library for logarithmic forms.

\section{Further deductions and refinements}
\label{sec:further}

\subsection{The first mismatch controls the direct defect}

Combining \cref{cor:first-mismatch-depth,lem:lead-lower} gives the fully
explicit estimate
\begin{equation}\label{eq:D-explicit-u}
 D_k\ge
 \begin{cases}
 (3^{L_k}+1)/2,
   &u_k\ge\log_3 2,\\[3pt]
 (3^{L_k-t(u_k)}+1)/2,
   &0<u_k<\log_3 2.
 \end{cases}
\end{equation}
This bound arises from digit order, not from convexity.  It can be stronger
than the smooth envelope when the first mismatch digit is high.

\subsection{The only dangerous leading mismatch}

If the leading coefficient of \(P_k\) is \(2\), then
\(E_k\ge2^{J_k}+1\) by \cref{lem:lead-lower}.  Therefore any substantial
cancellation in \(E_k=P_k(2)\) requires a leading coefficient \(1\), which
means that the first mismatching digit pair is exactly
\begin{equation}\label{eq:unit-mismatch-types}
 (\delta_{k,J_k},\epsilon_{k,J_k})\in\{(1,0),(2,1)\}.
\end{equation}
The following digits must then begin with the inverse pattern \((0,1)\) if
\(E_k\) is small relative to \(2^{J_k}\).  Hence every difficult word has
the schematic form
\begin{equation}\label{eq:word-schematic}
 \underbrace{(1,1)(0,0)\cdots(0,0)}_{\text{forced common top}}
 \quad
 \underbrace{(1,0)\text{ or }(2,1)}_{\text{positive mismatch}}
 \quad
 \underbrace{(0,1)\cdots(0,1)}_{\text{borrow block}}
 \quad\cdots.
\end{equation}
This is a much smaller symbolic language than the unrestricted product of a
binary and a ternary shift.

\subsection{A compactness reduction for symbolic attacks}

Fix a constant \(C\) large enough for the logarithmic bounds in
\cref{thm:prefix-log,cor:borrow-log}.  At exponent \(k\), all nontrivial
most-significant cancellation is confined to two windows of total length
\(O(C\log k)\): the common-zero window above \(J_k\) and any borrow window
below it.  Outside these windows, \(P_k(2)\) and \(P_k(3)\) have a stable
sign contribution of exponential size.

This suggests encoding only the \(O(\log k)\)-window as a state.  The number
of states is polynomial in \(k\), not exponential in the full digit length.
A successful multiplication-by-\(M\), multiplication-by-\(A\) transition
analysis would therefore need to beat polynomial state growth.  This is a
concrete place where an S-unit theorem, a subspace theorem with moving
targets, or an entropy estimate might enter.

\subsection{Why four exponentials still appears}

The new results exploit integrality much more directly than a bare
algebraicity argument, but they do not create an algebraic relation between
\(\log2\) and \(\log3\).  The synchronized ladder reduces two exponential
approximations to one continued-fraction error, yet the exact digit maps are
nonlinear and nonmultiplicative.  A closing theorem must therefore add
arithmetic rigidity of digital carries or prime factors.  Without that
input, the construction remains compatible with the known
four-exponentials frontier described in the supplied report.

\section{Conclusions}
\label{sec:conclusion}

The investigation produced the following rigorous advances.

\begin{enumerate}
\item Two exact digit-transfer maps \(\Phi,\Psi\) convert the hypothetical
power relation into positive integral defects.  Their carry algebra,
valuation increments, p-adic extensions, and sharp archimedean bounds are
fully explicit.

\item The same digit-difference polynomial evaluates to \(D_k\) at \(3\)
and \(E_k\) at \(2\).  Its leading coefficient is positive, and any strong
cancellation at \(2\) forces a long, unique \((0,1)\) borrow block.

\item The most-significant binary and ternary words cannot share an
arbitrary prefix.  Their common prefix is only \(1\) followed by zeros, and
the first mismatch favors the ternary digit.  This mismatch lies within
\(O(\log k)\) of the top for every \(k\), effectively for a fixed
hypothetical pair.

\item Both defects are macroscopic on a set of exponents of density
arbitrarily close to one.  An explicit logarithmic-form estimate strengthens
this to exponential lower bounds with only polynomial loss for every
exponent.

\item The pure-radical reduction yields norm gaps whose adjacent determinant
is bounded.  Optimizing the approximated logarithm produces a new ladder
with determinant exactly one and sharp adjacent coprimality.

\item On one continued-fraction sign subsequence, the ladder fits into an
exact gap-descent square.  A rigorous asymptotic calculation shows that the
square is noncontractive, preventing an illusory descent proof and exposing
the missing scale.
\end{enumerate}

The conjecture is therefore not proved here.  The most credible next step is
not another generic approximation to \(\vartheta\), but one of two highly
specific inputs: a prime-transfer theorem across the determinant-one
square, or an endpoint-sensitive valuation theorem for the exact defect
intervals.  Both targets now come with rigid digit words, logarithmic state
windows, and explicit failure thresholds, making them substantially more
falsifiable than the broad approaches they replace.

\begin{thebibliography}{99}

\bibitem{AlaogluErdos1944}
L. Alaoglu and P. Erd\H{o}s,
\newblock On highly composite and similar numbers,
\newblock \emph{Transactions of the American Mathematical Society}
\textbf{56} (1944), 448--469.

\bibitem{AlloucheShallit}
J.-P. Allouche and J. Shallit,
\newblock \emph{Automatic Sequences: Theory, Applications, Generalizations},
\newblock Cambridge University Press, 2003.

\bibitem{BakerWustholz}
A. Baker and G. W\"ustholz,
\newblock \emph{Logarithmic Forms and Diophantine Geometry},
\newblock Cambridge University Press, 2007.

\bibitem{Evertse1984}
J.-H. Evertse,
\newblock On equations in S-units and the Thue--Mahler equation,
\newblock \emph{Inventiones Mathematicae} \textbf{75} (1984), 561--584.

\bibitem{KuipersNiederreiter}
L. Kuipers and H. Niederreiter,
\newblock \emph{Uniform Distribution of Sequences},
\newblock Wiley, 1974; reprint, Dover, 2006.

\bibitem{Matveev2000}
E. M. Matveev,
\newblock An explicit lower bound for a homogeneous rational linear form in
logarithms of algebraic numbers. II,
\newblock \emph{Izvestiya: Mathematics} \textbf{64} (2000), 1217--1269.

\bibitem{SengeStraus}
H. G. Senge and E. G. Straus,
\newblock PV-numbers and sets of multiplicity,
\newblock \emph{Periodica Mathematica Hungarica} \textbf{3} (1973), 93--100.

\bibitem{Stewart1980}
C. L. Stewart,
\newblock On the representation of an integer in two different bases,
\newblock \emph{Journal f\"ur die reine und angewandte Mathematik}
\textbf{319} (1980), 63--72.

\bibitem{SuppliedReport}
\newblock \emph{Unified working report on the Alaoglu--Erd\H{o}s conjecture},
\newblock supplied internal manuscript, version A, 2026.

\end{thebibliography}

\end{document}
